\documentclass[11pt,a4paper]{article}
\usepackage[margin=2.5cm]{geometry}
\usepackage{amsmath,amssymb,amsthm,mathtools,enumitem}
\usepackage[colorlinks=true,linkcolor=blue!50!black]{hyperref}
\usepackage{booktabs}
\usepackage{xcolor}

\newtheorem{theorem}{Theorem}
\newtheorem{lemma}[theorem]{Lemma}
\newtheorem{proposition}[theorem]{Proposition}
\newtheorem{corollary}[theorem]{Corollary}
\theoremstyle{definition}
\newtheorem{definition}[theorem]{Definition}
\newtheorem{assumption}{Assumption}
\renewcommand{\theassumption}{A\arabic{assumption}}
\theoremstyle{remark}
\newtheorem{remark}[theorem]{Remark}
\newtheorem*{correction}{Correction to the draft}

\newcommand{\R}{\mathbb{R}}
\newcommand{\B}{\mathcal{B}}
\newcommand{\C}{\mathcal{C}}
\newcommand{\T}{\mathcal{T}}
\newcommand{\simp}{\Delta^{\circ}}
\newcommand{\eps}{\epsilon}
\newcommand{\RZ}{R_Z}
\newcommand{\ip}[2]{\langle #1,#2\rangle}

\title{Lemma S1$'$ and Lemma B: rigorous proofs\\[2pt]
\large with an audit of the review draft (E006, commit 4867fe3)}
\date{2026-11-21}
\author{}

\begin{document}
\maketitle

\paragraph{Scope.} This note restates the setting of the review draft with every
standing assumption made explicit, and proves Lemma~S1$'$ (items (1)--(4)), the
Gate~1A corollary, and Lemma~B. Where the draft's argument is correct, it is
reproduced with the missing details filled in. Where it is incomplete or
inaccurate, the change is marked \emph{Correction to the draft}. Facts about the
frozen specification (§5.3, §7, §9, the $9\times 9$ grid) are taken from the draft
as given; they are not mathematical statements and are not proved here.

\paragraph{Audit summary.}
\begin{center}
\begin{tabular}{@{}lll@{}}
\toprule
Item & Verdict & What changes \\
\midrule
Frame covariance & correct & assumption \ref{A:frame} made explicit; one-line proof added\\
S1$'$(1) N-q local & correct & finiteness needed; sharp identified set added (Prop.~\ref{prop:sharp})\\
S1$'$(2) N-loc global & correct & none\\
S1$'$(3) N-free global & correct & $1+\eps'\neq 0$ unnecessary; translation construction\\
& & keeps regularity (needed for consistency with (4))\\
S1$'$(4) zero eff.\ info & correct, gap & explicit in-class nuisance submodels; N-free needs\\
& & a positivity-preserving bump\\
Corollary (Gate 1A) & correct, gap & ``N-free is larger than N-q'' is false; use (3) for\\
& & N-free; well-specification not needed\\
Lemma B & correct & single-occurrence assumption made explicit;\\
& & smoothness for FTC; constant $\tfrac12$ is sharp\\
\bottomrule
\end{tabular}
\end{center}

%=====================================================================
\section{Preliminaries}

\subsection{Identification notions}
Let $\{P_{\eps,\eta}\}$ be a model with parameter of interest $\eps$ and nuisance
$\eta$, true point $(\eps_0,\eta_0)$. Write $(\eps,\eta)\sim(\eps',\eta')$
(\emph{observationally equivalent}) if $P_{\eps,\eta}=P_{\eps',\eta'}$. The
\emph{identified set} is $I(\eps_0,\eta_0):=\{\eps' : \exists\eta',\
(\eps',\eta')\sim(\eps_0,\eta_0)\}$. Local non-identification means
$I\cap(\eps_0-r,\eps_0+r)\neq\{\eps_0\}$ for all $r>0$; global
non-identification means $I$ equals the whole parameter set of $\eps$.

\begin{remark}[Zero information does not imply non-identification]
The draft's warning after Definition~1 is correct. Example: $P_\eps=N(\eps^3,1)$,
$\eps_0=0$. The score $3\eps^2 x$ vanishes at $0$, so the Fisher information is $0$,
yet $\eps\mapsto\eps^3$ is injective and $\eps$ is globally identified. This is why
every non-identification claim below is proved by construction.
\end{remark}

\subsection{Data and frame covariance}
Fix a finite candidate set $\C$ and the finite outcome set $\B=\{0,1\}^n$.
Candidate $c$ has tagged logical angles $\theta_c=(\theta_{c,1},\dots,\theta_{c,L_c})$
and a finite set $\T_c$ of label patterns with sign map $s:\T_c\to\{\pm1\}^{L_c}$,
drawn with known probabilities $\pi_c(t)>0$ and recorded. The candidate mix is
known (by design). The observed law is
$Q(c,t,b)=w_c\,\pi_c(t)\,P_{c,t}(b)$ with known $w_c,\pi_c$. Hence \emph{two
parameter points are observationally equivalent iff they induce the same family
$\{P_{c,t}\}_{c\in\C,\,t\in\T_c}$.} Write $\simp$ for the open probability simplex
on $\B$ and $\Sigma_0:=\{v\in\R^{\B}:\sum_b v(b)=0\}$.

\begin{assumption}[Frame covariance]\label{A:frame}
For tagged rotation $\ell$, the label is a Pauli $\sigma\in\{I,X,Y,Z\}$ on the
rotation's qubit, applied without error (or tracked virtually) on both sides;
$s_\ell=+1$ for $\sigma\in\{I,Z\}$ and $s_\ell=-1$ for $\sigma\in\{X,Y\}$. The
commanded angle is $a=s_\ell\theta_{c,\ell}$ and the physical gate is
$\RZ((1+\eps)a+\phi)$, $\RZ(x):=e^{-ixZ/2}$. Nothing else in the circuit or the
readout depends on $t$ except through these rotations.
\end{assumption}

\begin{lemma}[Frame covariance]\label{lem:frame}
Under \ref{A:frame} the label-$t$ circuit implements, exactly, the unlabelled
circuit with logical angles
$u_{c,\ell}(t)=(1+\eps)\theta_{c,\ell}+s_\ell(t)\phi$. Consequently
$P_{c,t}=R_c(u_c(t))$, where $R_c(u)$ is the outcome law of candidate $c$ with
logical angles $u$. At $\phi=0$, $P_{c,t}=R_c((1+\eps)\theta_c)=:P_c$ for every $t$.
\end{lemma}
\begin{proof}
$XZX=YZY=-Z$ and $IZI=ZZZ=Z$, so $\sigma Z\sigma=s Z$ and hence
$\sigma\RZ(x)\sigma=e^{-isxZ/2}=\RZ(sx)$ exactly (no global phase). Then
$\sigma\RZ((1+\eps)s\theta+\phi)\sigma=\RZ(s^2(1+\eps)\theta+s\phi)
=\RZ((1+\eps)\theta+s\phi)$ because $s^2=1$.
\end{proof}

\subsection{The three nuisance classes}
\begin{description}[leftmargin=1.2em]
\item[N-free.] $\eta=(R_c)_{c\in\C}$ with $R_c:\R^{L_c}\to\simp$ in a class
$\mathcal F_c$. The draft takes $\mathcal F_c$ = all maps. The proofs below use
only: \textup{(F1)} $\mathcal F_c$ is closed under translations
$R\mapsto R(\cdot+a)$; \textup{(F2)} for $R\in\mathcal F_c$, $u_0\in\R^{L_c}$,
$w\in\Sigma_0$, and a bump $\psi$ (below), $R+\tau\psi w\in\mathcal F_c$ whenever it
maps into $\simp$. All maps, continuous maps and $C^k$ maps all satisfy (F1)--(F2).
\item[N-loc.] $\eta=(q_c,D_{c,1},\dots,D_{c,L_c})_c$ with $q_c\in\simp$,
$D_{c,\ell}\in\Sigma_0$ and
$P_{c,t}=q_c+\sum_\ell(\eps\theta_{c,\ell}+s_\ell(t)\phi)D_{c,\ell}$,
subject to $P_{c,t}\in\simp$ for all $t\in\T_c$.
\item[N-q.] As N-loc with $D=D^0$ known; $\eta=(q_c)_c$.
\end{description}
Set $v_c:=\sum_\ell\theta_{c,\ell}D_{c,\ell}\in\Sigma_0$. At $\phi=0$:
$P_{c,t}=q_c+\eps v_c$ for all $t$ (N-loc, N-q).

\begin{correction}
The draft (Corollary, Step 4) calls N-loc and N-free ``larger classes'' than N-q.
N-q $\subseteq$ N-loc is true. N-q $\subseteq$ N-free is \emph{false}: the
N-loc law is $R_c(u)=q_c+\sum_\ell(u_\ell-\theta_{c,\ell})D_{c,\ell}$, which is affine
in $u$; if some $D_{c,\ell}\neq0$ it leaves $\simp$ for large $|u|$, so it is not a
map $\R^{L_c}\to\simp$. The classes N-loc and N-free are not nested. This does not
affect any conclusion, because item (3) handles N-free directly
(see Corollary~\ref{cor:gate}).
\end{correction}

%=====================================================================
\section{Lemma S1$'$}

\begin{theorem}[Lemma S1$'$]\label{thm:S1}
Fix a true point with $\phi_0=0$ (so $P_{c,t}=P_c\in\simp$ for all $t$).
\begin{enumerate}[label=(\arabic*)]
\item Under N-q, and hence under N-loc, $I\supseteq(\eps_0-\delta,\eps_0+\delta)$
for some $\delta\in(0,\infty]$.
\item Under N-loc, $I=\R$.
\item Under N-free (with (F1)), $I=\R$.
\item If, under N-free, $R_c$ is differentiable at $u^0_c:=(1+\eps_0)\theta_c$
for each $c$, then under each of N-q, N-loc, N-free (with (F2)) the score for $\eps$
lies in the nuisance tangent set itself. Hence the efficient score is $0$ and the
efficient information for $\eps$ is $0$.
\end{enumerate}
\end{theorem}

\subsection{Proof of (1)}
Let $q^0_c$ be the true intercepts, $m:=\min_{c\in\C}\min_{b\in\B}q^0_c(b)$, and
$V:=\max_c\|v_c\|_\infty$. Both are finite minima/maxima over finite sets, and
$m>0$ because each $q^0_c\in\simp$. Put $\delta:=m/V$ ($\delta:=\infty$ if $V=0$).

Fix $\eps'$ with $|\eps'-\eps_0|<\delta$ and define $q'_c:=q^0_c+(\eps_0-\eps')v_c$,
$D':=D^0$, $\phi':=0$.
\emph{Sum:} $\sum_b q'_c(b)=1+(\eps_0-\eps')\sum_bv_c(b)=1$.
\emph{Positivity:} if $v_c(b)=0$ then $q'_c(b)=q^0_c(b)>0$; otherwise $V>0$ and
$|(\eps_0-\eps')v_c(b)|<\delta V=m\le q^0_c(b)$. So $q'_c\in\simp$.
\emph{Laws:} for every $t$, $P'_{c,t}=q'_c+\eps'v_c=q^0_c+\eps_0v_c=P_c\in\simp$,
so the N-loc/N-q positivity constraint holds. Thus
$(\eps',0,q')\sim(\eps_0,0,q^0)$ under N-q. The same point, with $D'=D^0$, lies in
the N-loc parameter space, so the conclusion transfers to N-loc. \qed

\begin{proposition}[Sharp identified set under N-q]\label{prop:sharp}
Under N-q, suppose for some $c$ the map
$t\mapsto\sum_\ell s_\ell(t)D^0_{c,\ell}$ is non-constant on $\T_c$ (or $\phi$ is known
to be $0$). Then
\[
I=\bigl(\eps_0-\alpha,\ \eps_0+\beta\bigr),\quad
\alpha=\min_{c,b:\,v_c(b)<0}\frac{q^0_c(b)}{|v_c(b)|},\quad
\beta=\min_{c,b:\,v_c(b)>0}\frac{q^0_c(b)}{v_c(b)}
\]
(minimum over an empty set $=\infty$). If some $v_c\neq0$, then $\alpha,\beta<\infty$,
so $\eps$ is \emph{locally but not globally} non-identified.
\end{proposition}
\begin{proof}
Let $(\eps',\phi',q')\sim$ truth. For the $c$ in the hypothesis,
$P'_{c,t}-P'_{c,t'}=\phi'\sum_\ell(s_\ell(t)-s_\ell(t'))D^0_{c,\ell}$ must vanish for all
$t,t'$, forcing $\phi'=0$. Then $q'_c+\eps'v_c=q^0_c+\eps_0v_c$ for every $c$, i.e.\
$q'_c$ is determined, and the only constraint left is
$q^0_c(b)-(\eps'-\eps_0)v_c(b)>0$ for all $c,b$, which is exactly
$-\alpha<\eps'-\eps_0<\beta$. Since $v_c\in\Sigma_0$, $v_c\neq0$ has both a positive and
a negative entry, so both minima are over non-empty sets.
\end{proof}
This makes precise the draft's §2.6 remark that under N-q ``only local
non-identification is guaranteed'': under the stated condition, global
non-identification actually fails. Note $\delta=m/V\le\min(\alpha,\beta)$, consistent
with (1).

\subsection{Proof of (2)}
Let $\eps'\in\R$; set $\phi'=0$.
If $\eps'\neq0$: $D'_{c,\ell}:=(\eps_0/\eps')D^0_{c,\ell}\in\Sigma_0$, $q'_c:=q^0_c$.
Then $\eps'v'_c=\eps_0v_c$ and $P'_{c,t}=q^0_c+\eps_0v_c=P_c$ for all $t$.
If $\eps'=0$: $q'_c:=P_c\in\simp$, $D':=0$; then $P'_{c,t}=P_c$.
In both cases $q'_c\in\simp$ and $P'_{c,t}=P_c\in\simp$, so the point is admissible
and equivalent to the truth. \qed

\subsection{Proof of (3)}
Let $\eps'\in\R$ and define $R'_c(u):=R^0_c\bigl(u+(\eps_0-\eps')\theta_c\bigr)$,
$\phi':=0$. By (F1), $R'_c\in\mathcal F_c$; it maps into $\simp$ and has exactly the
regularity of $R^0_c$. By Lemma~\ref{lem:frame}, for every $t$,
$P'_{c,t}=R'_c((1+\eps')\theta_c)=R^0_c((1+\eps_0)\theta_c)=P_c$. \qed

\begin{correction}
(i) The restriction $1+\eps'\neq0$ in the draft is unnecessary: the construction
works for every real $\eps'$ (also the draft's pointwise one: if $\eps'=-1$ the single
observed point is $u=0$, and $R'_c(0)$ may still be set to $P_c$).
(ii) The draft's pointwise construction (``any strictly positive distribution
elsewhere'') is valid only if $\mathcal F_c$ contains arbitrary maps. If the frozen
class requires continuity or smoothness---which item (4) implicitly does---the
pointwise $R'_c$ may leave the class. The translation above is valid for every class
closed under translations, so it covers all readings.
\end{correction}

\subsection{Proof of (4)}
Work at the truth $(\eps_0,\phi_0{=}0,\eta_0)$. Because $w_c,\pi_c$ are known, the
log-likelihood of one observation is $\log P_{c,t}(b)$ plus a known constant, and every
score has the form $h(c,t,b)=\dot P_{c,t}(b)/P_{c,t}(b)$, where $\dot P$ is the
derivative of the law along a one-dimensional submodel. All $P_{c,t}(b)>0$, so these are
bounded functions in $L^2(Q)$.

\emph{Score for $\eps$.} N-loc/N-q: $\partial_\eps P_{c,t}=v_c$ (for all $\phi$),
so $\dot\ell_\eps=v_c(b)/P_c(b)$. N-free: by the chain rule at $\phi=0$,
$\partial_\eps R_c((1+\eps)\theta_c+s(t)\phi)\big|_{\eps_0,0}
=\nabla R_c(u^0_c)\theta_c=:w_c$, so $\dot\ell_\eps=w_c(b)/P_c(b)$. In both cases the
score does not depend on $t$, and $v_c,w_c\in\Sigma_0$ (for $w_c$: differentiate
$\sum_bR_c(u)(b)\equiv1$ along $\theta_c$).

\emph{A nuisance submodel with the same score.}
\begin{itemize}[leftmargin=1.2em]
\item N-q and N-loc: $q^\tau_c:=q^0_c+\tau v_c$, all else fixed. For
$|\tau|<m/V$, $q^\tau_c\in\simp$ and $P^\tau_{c,t}=P_c+\tau v_c\in\simp$, so the path
stays in the class (in N-q, $D$ is not touched). Its score at $\tau=0$ is
$v_c(b)/P_c(b)=\dot\ell_\eps$.
\item N-free: $R_c$ is continuous at $u^0_c$ (differentiability implies continuity), so
there is $r>0$ with $\min_bR_c(u)(b)\ge m_c/2$ for $\|u-u^0_c\|<r$, where
$m_c:=\min_bP_c(b)>0$. Let $\psi_c:\R^{L_c}\to[0,1]$ be a $C^\infty$ bump with
$\psi_c(u^0_c)=1$ and support in that ball. Put $R^\tau_c:=R_c+\tau\psi_cw_c$. Each
value sums to $1$; outside the ball it equals $R_c$; inside it is $\ge m_c/2-|\tau|\,
\|w_c\|_\infty>0$ for $|\tau|<m_c/(2\|w_c\|_\infty)$. By (F2) the path lies in the
class, and its score at $\tau=0$ is $\psi_c(u^0_c)w_c(b)/P_c(b)=\dot\ell_\eps$.
\end{itemize}
Hence $\dot\ell_\eps\in\dot{\mathcal P}_\eta$ (the nuisance tangent set, not merely its
closed span). The efficient score is
$\tilde\ell_\eps=\dot\ell_\eps-\Pi(\dot\ell_\eps\mid\overline{\mathrm{lin}}\,
\dot{\mathcal P}_\eta)=0$, and the efficient information
$\tilde I=E_Q[\tilde\ell_\eps^{\,2}]=0$. Enlarging the nuisance (e.g.\ treating
$\phi$ as unknown) only enlarges $\dot{\mathcal P}_\eta$, so the conclusion is unchanged.
\qed

\begin{remark}
The argument shows more than zero information: the direction
$\eps\mapsto\eps+\tau$ coincides, to first order, with a nuisance direction. Item (1)
shows the coincidence is exact (not only first-order) under N-q, which is why
non-identification and zero information hold together here, unlike the
$N(\eps^3,1)$ example.
\end{remark}

%=====================================================================
\section{Gate 1A corollary}

Let $\Theta$ be the frozen class (any one of N-q, N-loc, N-free). At exact
probabilities the population log-likelihood criterion is
\begin{gather*}
\Lambda(\eps',\phi',\eta')=\sum_cw_c\sum_t\pi_c(t)\sum_bP_{c,t}(b)\log P'_{c,t}(b),\\
\Lambda^*:=\sum_cw_c\sum_t\pi_c(t)\sum_bP_{c,t}(b)\log P_{c,t}(b),
\end{gather*}
and the profile is $\Lambda_{\rm pr}(\eps'):=\sup_{\phi',\eta'}\Lambda(\eps',\phi',\eta')$.

\begin{lemma}[Gibbs]\label{lem:gibbs}
$\Lambda\le\Lambda^*$, with equality iff $P'_{c,t}=P_{c,t}$ for all $(c,t)$.
\end{lemma}
\begin{proof}
$\Lambda^*-\Lambda=\sum_cw_c\sum_t\pi_c(t)\,\mathrm{KL}(P_{c,t}\|P'_{c,t})\ge0$, with
equality iff every KL term is $0$, since $w_c,\pi_c(t)>0$.
\end{proof}

\begin{corollary}\label{cor:gate}
At any true point with $\phi_0=0$, the set of maximisers of $\Lambda_{\rm pr}$ that
attain $\Lambda^*$ equals the identified set $I$. It contains an open interval
around $\eps_0$ under N-q and N-loc (Theorem~\ref{thm:S1}(1)) and equals $\R$ under
N-loc and N-free (Theorem~\ref{thm:S1}(2),(3)). In every reading the profile
maximiser is not unique.
\end{corollary}
\begin{proof}
By Lemma~\ref{lem:gibbs}, $\Lambda_{\rm pr}(\eps')=\Lambda^*$ is attained iff there is
$(\phi',\eta')$ with $P'=P$, i.e.\ iff $\eps'\in I$. Then apply Theorem~\ref{thm:S1}.
\end{proof}

\begin{remark}[Well-specification is not needed]
Suppose only that the true per-candidate laws are label-invariant and in $\simp$, with
no claim that they are of N-q form. Under N-q, for every $\eps'$ with
$P_c-\eps'v_c\in\simp$ for all $c$ (an open interval containing $0$), the choice
$q'_c=P_c-\eps'v_c$, $\phi'=0$ gives $P'=P$, so $\Lambda_{\rm pr}=\Lambda^*$ there. The
profile is flat on an interval even if the class is misspecified.
\end{remark}

\paragraph{From non-uniqueness to FAIL (specification step, not mathematics).}
Given the draft's description of the frozen rule---§7: profile ML under the I3 class,
applied to exact probabilities, with a flat profile or a nonunique solution being a
failure; §9: recovery required at every grid point; nine grid points with $\phi=0$---the
corollary implies FAIL at each of those nine points, under every reading. This step
uses the specification's wording, which should be checked against the repository text.
It does not depend on what any particular optimiser returns from its starting points.

%=====================================================================
\section{Lemma B}

\subsection{Setting and assumptions}
\begin{assumption}[Single occurrence]\label{A:single}
For $j\in\{A,B\}$,
\[U_j(x_j)=W_{j,m_j}\,\RZ^{(k_{m_j})}(x_{j,m_j})\,W_{j,m_j-1}\cdots
W_{j,1}\,\RZ^{(k_1)}(x_{j,1})\,W_{j,0},\]
where the $W_{j,i}$ are fixed unitaries
independent of all tagged variables (CX gates, mixers, untagged rotations), and
$\RZ^{(k)}$ acts on qubit $k$. Each tagged variable occurs in exactly one gate.
(If a QAOA angle $\gamma$ appears in several gates, each occurrence is a separate
coordinate $x_{j,i}$, with its own perturbation $\delta_{j,i}$, as in the error model.)
\end{assumption}
\begin{assumption}[Error model, no labels]\label{A:err}
Target circuits carry no RC labels, so $s=+1$ and the commanded angle is $\theta_{j,i}$.
The physical gate is $\RZ((1+\eps)\theta_{j,i}+\phi)=\RZ(\theta_{j,i}+\delta_{j,i})$
with $\delta_{j,i}=\eps\theta_{j,i}+\phi$, exactly. All other gates are error-free.
\end{assumption}
Let $s:\B\to[0,1]$,
$E_j(x_j)=\sum_bs(b)\,|\langle b|U_j(x_j)|{+}^n\rangle|^2$,
$g(x)=E_A(x_A)-E_B(x_B)$, $A=\sum_{j,i}|\theta_{j,i}|$, $M=m_A+m_B$,
$B(\eps,\phi)=|\eps|A+|\phi|M$. (The draft's condition $\theta_{j,i}>0$ is not needed.)

\begin{theorem}[Lemma B]\label{thm:B}
For all $\eps,\phi\in\R$,
\[|g(\theta+\delta)-g(\theta)|\le\tfrac12\|\delta\|_1\le\tfrac12B(\eps,\phi).\]
\end{theorem}

\subsection{Proof}
\paragraph{Step 0 (smoothness).} Each matrix entry of $\RZ(x)$ is $\cos(x/2)$ or
$\pm i\sin(x/2)$ times a constant. Hence each amplitude, and so $E_j$, is a polynomial in
$\{\cos(x_{j,i}/2),\sin(x_{j,i}/2)\}_i$; in particular $E_j\in C^\infty(\R^{m_j})$ and
$g\in C^\infty(\R^{M})$.

\paragraph{Step 1 (one angle).} Fix all coordinates except $x=x_{j,i}$, at arbitrary
values. Then $U_j=V_2\RZ^{(k)}(x)V_1$ with $V_1,V_2$ independent of $x$, and
$\RZ^{(k)}(x)=\cos\tfrac x2\,I-i\sin\tfrac x2\,Z_k$. So
$\langle b|U_j|{+}^n\rangle=\cos\tfrac x2\,\alpha_b-i\sin\tfrac x2\,\beta_b$ with
$\alpha_b=\langle b|V_2V_1|{+}^n\rangle$, $\beta_b=\langle b|V_2Z_kV_1|{+}^n\rangle$.
Expanding,
\[
|\cdot|^2=\cos^2\tfrac x2|\alpha_b|^2+\sin^2\tfrac x2|\beta_b|^2
-2\sin\tfrac x2\cos\tfrac x2\,\mathrm{Im}(\alpha_b\overline{\beta_b}),
\]
and with $\cos^2\tfrac x2=\tfrac{1+\cos x}2$, $\sin^2\tfrac x2=\tfrac{1-\cos x}2$,
$2\sin\tfrac x2\cos\tfrac x2=\sin x$:
\begin{gather*}
E_j=a+b\cos x+c\sin x,\\
a=\tfrac12\sum_bs(b)(|\alpha_b|^2+|\beta_b|^2),\quad
b=\tfrac12\sum_bs(b)(|\alpha_b|^2-|\beta_b|^2),\\
c=-\sum_bs(b)\,\mathrm{Im}(\alpha_b\overline{\beta_b}).
\end{gather*}

\paragraph{Step 2 (uniform derivative bound).}
\begin{lemma}\label{lem:trig}
If $f(x)=a+b\cos x+c\sin x$ satisfies $0\le f\le1$ on $\R$, then $|f'(x)|\le\tfrac12$
for all $x$.
\end{lemma}
\begin{proof}
Let $\rho=\sqrt{b^2+c^2}$. If $\rho=0$, $f'\equiv0$. Otherwise write
$f=a+\rho\cos(x-x_0)$; evaluating at $x=x_0$ and $x=x_0+\pi$ gives $a+\rho\le1$ and
$a-\rho\ge0$, so $\rho\le\tfrac12$. Then $|f'(x)|=\rho|\sin(x-x_0)|\le\tfrac12$.
\end{proof}
For every value of $x$ (every unitary $U_j$), $E_j=\sum_bs(b)p(b)$ with $p$ a
probability vector and $s\in[0,1]$, so $0\le E_j\le1$. By Step 1 and
Lemma~\ref{lem:trig}, $|\partial E_j/\partial x_{j,i}|\le\tfrac12$ at every point of
$\R^{m_j}$ (the fixed coordinates in Step 1 were arbitrary).

\paragraph{Step 3 (global bound).} Since $x_A$ and $x_B$ are disjoint coordinates,
$\partial g/\partial x_{A,i}=\partial E_A/\partial x_{A,i}$ and
$\partial g/\partial x_{B,i}=-\partial E_B/\partial x_{B,i}$, so
$\|\nabla g(y)\|_\infty\le\tfrac12$ for all $y\in\R^M$. As $g\in C^1$, the fundamental
theorem of calculus along $y(t)=\theta+t\delta$ gives
\[
|g(\theta+\delta)-g(\theta)|=\Bigl|\int_0^1\nabla g(y(t))\cdot\delta\,dt\Bigr|
\le\int_0^1\|\nabla g(y(t))\|_\infty\|\delta\|_1\,dt\le\tfrac12\|\delta\|_1,
\]
using Hölder's inequality $|\ip{u}{v}|\le\|u\|_\infty\|v\|_1$.

\paragraph{Step 4 (size of $\delta$).} By \ref{A:err},
$|\delta_{j,i}|=|\eps\theta_{j,i}+\phi|\le|\eps||\theta_{j,i}|+|\phi|$; summing over
the $M$ tagged rotations gives $\|\delta\|_1\le|\eps|A+|\phi|M=B(\eps,\phi)$. \qed

\subsection{Remarks}
\begin{itemize}[leftmargin=1.2em]
\item \emph{Sharpness of $\tfrac12$.} Over the class of circuits and scores allowed by
\ref{A:single}, the constant cannot be lowered: one qubit, $W_0=I$, $W_1=H$,
$s(0)=1$, $s(1)=0$ gives $E(x)=\cos^2(x/2)$, $|E'(\pi/2)|=\tfrac12$, and
$|\Delta g|/(\tfrac12\|\delta\|_1)\to1$ as $\delta\to0$ along that coordinate. For the
frozen QAOA instances the observed ratio ($\le0.431$) is smaller, which is an
instance property, not a contradiction.
\item \emph{A trivial companion bound.} Since $E_j\in[0,1]$,
$|\Delta g|\le\min(\tfrac12\|\delta_A\|_1,1)+\min(\tfrac12\|\delta_B\|_1,1)$, which is
better for large perturbations.
\item \emph{Wording.} The draft's ``$B$ is exactly twice a global bound'' is best stated
as: $B/2$ is a global, exact (no remainder) bound on $|\Delta g|$; hence $B$ is a bound
with a factor-$2$ margin. This correctly refutes E003~§6.
\item \emph{What is not proved.} The draft's §3.6 list is accurate and should stand:
nothing here concerns $\hat B=B(\hat\eps,\hat\phi)$, coverage, shot noise in $\hat g_0$,
decision-flip probabilities, post-recalibration residuals, or models beyond the
noiseless synthetic E2 model with errors on tagged rotations only.
\end{itemize}

%=====================================================================
\section{Suggested checklist outcomes}
\begin{center}
\begin{tabular}{@{}p{0.55\textwidth}p{0.38\textwidth}@{}}
\toprule
Draft checklist item & Suggested outcome\\
\midrule
Frame covariance $\Rightarrow$ label invariance & accept; cite \ref{A:frame}\\
Three parameterizations cover §5.3 & spec question; not checkable here\\
(1) N-q construction & accept; add finiteness, Prop.~\ref{prop:sharp}\\
(2) N-loc construction & accept\\
(3) N-free construction & accept with correction (translation)\\
(4) score in nuisance tangent span & accept with correction (explicit paths)\\
Corollary $\Rightarrow$ FAIL & accept with correction (non-nesting)\\
``Does not claim'' list & accept\\
Lemma B Steps 1--4 & accept; add \ref{A:single}, Step 0\\
Unproved items about $\hat B$ & accept as stated\\
\bottomrule
\end{tabular}
\end{center}
Overall: \textbf{accept with corrections}. No correction changes a conclusion.

\end{document}
